% Figure: an AC impedance bridge (Maxwell-Wien form). Four impedance arms in
% a diamond, an AC source across one diagonal and a phase-sensitive detector
% across the other. Balance requires the complex cross-product equality
% Z1 Z4 = Z2 Z3, one equation in complex numbers, hence two real conditions.
\begin{figure}[htbp]
\centering
\begin{tikzpicture}[scale=1.0, line cap=round, >=latex]
  % four nodes of the diamond
  \coordinate (T) at (0, 2.2);   % top
  \coordinate (R) at (2.6, 0);   % right
  \coordinate (B) at (0, -2.2);  % bottom
  \coordinate (L) at (-2.6, 0);  % left
  % node dots
  \foreach \p in {T,R,B,L} \filldraw[engblack] (\p) circle (0.04);
  \node[font=\footnotesize, anchor=south] at (T) {$a$};
  \node[font=\footnotesize, anchor=west]  at (R) {$b$};
  \node[font=\footnotesize, anchor=north] at (B) {$c$};
  \node[font=\footnotesize, anchor=east]  at (L) {$d$};
  % arms as labelled boxes on the four sides
  \draw[engblue, very thick] (T) -- node[midway, above right, font=\footnotesize] {$Z_1$} (R);
  \draw[engblue, very thick] (R) -- node[midway, below right, font=\footnotesize] {$Z_4$} (B);
  \draw[engblue, very thick] (B) -- node[midway, below left, font=\footnotesize] {$Z_3$} (L);
  \draw[engblue, very thick] (L) -- node[midway, above left, font=\footnotesize] {$Z_2$} (T);
  % small rectangles on each arm to read as components
  \filldraw[white, draw=engblue, thick] ($(T)!0.5!(R)$) ++(-0.16,-0.11) rectangle ++(0.32,0.22);
  \filldraw[white, draw=engblue, thick] ($(R)!0.5!(B)$) ++(-0.16,-0.11) rectangle ++(0.32,0.22);
  \filldraw[white, draw=engblue, thick] ($(B)!0.5!(L)$) ++(-0.16,-0.11) rectangle ++(0.32,0.22);
  \filldraw[white, draw=engblue, thick] ($(L)!0.5!(T)$) ++(-0.16,-0.11) rectangle ++(0.32,0.22);
  % source across a-c diagonal (top to bottom)
  \draw[engred, thick] (T) .. controls (-1.1, 1.5) and (-1.1, -1.5) .. (B);
  \draw[engred, thick, fill=white] (-1.15, 0) circle (0.28);
  \node[engred, font=\small] at (-1.15, 0) {$\sim$};
  \node[engred, font=\footnotesize, anchor=east] at (-1.5, 0) {$\tilde V_s$};
  % detector across d-b diagonal (left to right)
  \draw[engamber, thick] (L) -- (R);
  \draw[engamber, thick, fill=white] (0, 0) circle (0.30);
  \node[engamber, font=\small] at (0, 0) {D};
  \node[engamber, font=\footnotesize, anchor=south] at (0, 0.34) {null detector};
\end{tikzpicture}
\caption[An AC impedance bridge]{An AC impedance bridge. Four impedance arms
$Z_1, Z_2, Z_3, Z_4$ form a diamond; a sinusoidal source $\tilde V_s$ drives
one diagonal ($a$ to $c$) and a phase-sensitive null detector D reads the
other ($d$ to $b$). The detector voltage is the difference of two voltage-
divider outputs, and it vanishes when the two dividers split the source in the
same complex ratio, the balance condition $Z_1 Z_4 = Z_2 Z_3$. That single
equation between complex numbers is two real equations at once, one for the
magnitude and one for the phase, so a working bridge carries two adjustable
elements and the operator nulls both the in-phase and the quadrature parts of
the detector reading. Balancing an AC bridge is the physical act of solving
one complex equation by hand.}
\label{fig:vol02ch03:ac-bridge}
\end{figure}
